\documentclass[11pt,a4paper]{article}
\usepackage[margin=27mm]{geometry}
\usepackage{amsmath,amssymb,amsthm,hyperref}
\newtheorem{lemma}{Lemma}\newtheorem{theorem}[lemma]{Theorem}
\newcommand{\Sz}{\operatorname{Sz}}\newcommand{\F}{\mathbb F}
\title{Scalar subspaces and unbounded algebraic local rank}
\author{Human Compression project}\date{}
\begin{document}\maketitle
\begin{abstract}
The algebraic residue convention in the fixed rank-three Suzuki article
admits locally doubly transitive, residually weakly primitive
pregeometries of arbitrarily large rank. A coordinate subspace supplies
the soluble parabolic; deleting its coordinates matches intersections
of independently chosen maximal subfields. This refutes the article's
separate universal rank assertion. No claim about full incidence
residues in higher rank is needed.
\end{abstract}

\section{The assertion and acquired structure}
The introduction of D.~Leemans, \emph{The rank 3 geometries of the simple
Suzuki groups $\Sz(q)$}, Note di Matematica 19(1) (1999), 43--63,
asserts on printed pp.~44--45 that no pregeometry of rank greater than
three satisfies both residual weak primitivity and local double
transitivity, and that rank three occurs only when $q-1$ is prime.
These also follow from the complete list stated in Theorem~6.1.
We use the definitions on pp.~45--46: the residue at a flag of the
distinguished chamber is formed from the parabolic intersections.
The local condition is imposed on these algebraic rank-one residues,
not on all elements in an actual incidence panel. This definition
applies to every finite type set; it is not restricted to rank three.
Flag transitivity is not assumed.

We retain the classical faithful Suzuki coordinates and subgroup theorem.
The standard group over a subfield consists of matrices with parameters
in that subfield. Its point stabilizer is $B_t=U_tH_t$, of order
$t^2(t-1)$. The standard involution centre is
$Z_F=\{u(0,b):b\in F\}$, with addition in its parameter. It centralizes
$U_t$, and $H_t$ acts by the full scalar group $\F_t^*$.
Indeed the standard action is $b\mapsto k^{\theta+2}b$ and
$\gcd(\theta+2,t-1)=1$.
These identities and the overgroup correspondence below are acquired
from the accepted scalar-subspace and subfield-family notes.
A simple subfield group is maximal at prime relative field degree.
The action of $\Sz(t)$ on $\Sz(t)/B_t$ is its classical doubly
transitive action on $t^2+1$ points.

\section{Independent field coordinates}
Let $t=2^e\ge8$, with $e$ odd. Choose distinct odd primes
$a_1,\ldots,a_k$, with $k\ge2$, and put
\[
 A=\prod_{i=1}^k a_i,\quad F=\F_{t^A},\quad K=\F_t,
 \quad F_i=\F_{t^{A/a_i}}.
\]
Choose $x_i\in\F_{t^{a_i}}\setminus K$, and set
\[
 T=K+Kx_1+\cdots+Kx_k,\qquad
 T_J=K+\sum_{i\notin J}Kx_i\quad(J\subseteq\{1,\ldots,k\}).
\]
\begin{lemma}\label{field}
The vectors $1,x_1,\ldots,x_k$ are linearly independent over $K$.
Moreover
\[
 F_J:=\bigcap_{j\in J}F_j=\F_{t^{A/\prod_{j\in J}a_j}},
 \qquad T\cap F_J=T_J,
\]
where the empty intersection means $F$.
\end{lemma}
\begin{proof}
The field intersection formula follows from the gcd of the degrees.
The automorphism $\sigma_i:y\mapsto y^{t^{A/a_i}}$ fixes $F_i$ and
all $x_j$ for $j\ne i$. Its restriction to $\F_{t^{a_i}}$ has fixed
field $K$, since $a_i$ is coprime to $A/a_i$. Hence
$\sigma_i(x_i)\ne x_i$. Applying $\sigma_i-1$ to a relation
$c_0+\sum c_jx_j=0$, with $c_j\in K$, gives
$c_i(\sigma_i(x_i)-x_i)=0$. Thus every $c_i=0$, and then $c_0=0$.
A vector in $T\cap F_i$ must likewise have zero $x_i$ coefficient.
Imposing this for all $i\in J$ proves the second equality.
\end{proof}

For a $K$-subspace $W\supseteq K$ of $F$, write
\[
 Q_W=Z_WB_t,\qquad Z_W=\{u(0,b):b\in W\}.
\]
This is a group: $U_t$ centralizes $Z_W$, and $H_t$ preserves $W$.
Its order is $t|W|(t-1)$, since $Z_W\cap B_t=Z_K$ has order $t$.
Every overgroup of $Q_V$ inside $Q_W$ is $Q_D$ for a $K$-subspace
$V\le D\le W$. Indeed every element has form $zb$, with
$z\in Z_W$, $b\in B_t$. A subgroup containing $B_t$ is therefore
of this form with an additive subgroup of $W$ containing $K$.
Normalization by $H_t$ makes that subgroup a $K$-subspace. Consequently
\[
 Q_V\text{ is maximal in }Q_W\quad\Longleftrightarrow\quad
 \dim_K(W/V)=1.                                      \tag{1}
\]
For any intermediate field $E\supseteq K$, unique coordinates give
\[
 \Sz(E)\cap Q_W=Q_{W\cap E}.                         \tag{2}
\]
In fact multiplying by an element of $B_t\le\Sz(E)$ reduces
membership to the central parameter in $W\cap E$.

\section{The construction and every residue check}
\begin{theorem}\label{family}
In $G=\Sz(t^A)$ take the $k+1$ parabolics
\[
 G_i=\Sz(F_i)\quad(1\le i\le k),\qquad G_*=Q_T.
\]
Their coset pregeometry satisfies the source's residual weak primitivity
and algebraic local double transitivity conditions. Its rank is $k+1$.
For $k=2$ it is an incidence geometry of rank three, and $q-1$ is
composite. Thus the source's general rank bound and its restriction
of rank three to prime $q-1$ are false.
\end{theorem}
\begin{proof}
For $J\subseteq\{1,\ldots,k\}$, unique field coordinates,
Lemma~\ref{field}, and (2) give all intersections:
\[
 G_J=\Sz(F_J),\qquad G_{J\cup\{*\}}=Q_{T_J}.          \tag{3}
\]
The intersection of all parabolics is $B_t=Q_K$. Every parabolic
is proper: each $G_i$ is a prime-degree subfield, and $Q_T$ is
contained in the ambient point stabilizer. The intersections in (3)
decrease strictly when any type is added.

We verify residual weak primitivity for every proper type subset.
If the subset omits $*$ and some $i$, adjoining $i$ is a prime-degree
simple subfield inclusion, hence maximal. If it contains all
$1,\ldots,k$ and omits only $*$, its group is $\Sz(t)$ and adjoining
$*$ gives its maximal point stabilizer $B_t$. If the subset contains
$*$ and omits an $i$, adjoining $i$ removes exactly the independent
coordinate $Kx_i$ from $T_J$. Equation~(1) gives maximality.
This exhausts every proper subset. The last index is $t$, not a prime;
maximality follows from scalar irreducibility as in (1).

The rank-one algebraic residue obtained by omitting $*$ is the
coset action of $\Sz(t)$ on $\Sz(t)/B_t$, doubly transitive by the
classical action. Omitting $i$ gives $Q_{K+Kx_i}$ on its cosets of
$B_t$. These cosets identify with $(K+Kx_i)/K\cong K$.
The centre provides all translations, while $H_t$ provides every
nonzero scalar multiplication. Its induced action is therefore the
affine group $K\rtimes K^*$, doubly transitive on $t$ points.
These are all the source's required local actions.

For $k=2$, every rank-three coset pregeometry is a geometry: an
incident pair has a common representative, which supplies the third
coset and extends it to a chamber. Singletons extend in the same way.
No flag transitivity follows from this observation. Finally
$q-1=(t-1)(1+t+\cdots+t^{A-1})$ has both factors greater than one.
The rank-three example thus has composite $q-1$. Since $k$ is
arbitrary, (3) gives pregeometries of arbitrarily large rank.
\end{proof}

Concrete parameters are $t=8$, $(a_1,a_2)=(3,5)$, giving rank
three at $q=2^{45}$; and $t=8$, $(a_1,a_2,a_3)=(3,5,7)$,
giving rank four at $q=2^{315}$. These are symbolic examples:
any $x_i$ in its displayed subfield outside $K$ suffices.
No ambient Suzuki group is enumerated.

\paragraph{Attribution and scope.}
Suzuki structure and its natural doubly transitive action are classical
premises. The correspondence $W\mapsto Q_W$ and its maximality criterion
reuse the accepted scalar-subspace theory. The contribution here is
the independent field-coordinate construction and its application to
the source's separate universal rank claim. No full incidence thinness,
higher-rank geometry property, or flag transitivity is asserted.
This result does not dispose of a different statement using full actual
panels or additional geometric hypotheses. It is an additional correction
within the already refuted original, not new manuscript or principal
refutation credit. No computation, Lean declaration or compression
factor is introduced.
\end{document}
